\documentclass[aspectratio=169,10pt]{beamer}

% Shared Beamer setup for Fall 2026 course slides.
% Clean Cal Poly-inspired theme: no custom class files or uncommon packages.
\mode<presentation>{
  \usetheme{default}
  \usefonttheme{professionalfonts}
}

\definecolor{CalPolyGreen}{HTML}{154734}
\definecolor{CalPolyGold}{HTML}{BD8B13}
\definecolor{CalPolySage}{HTML}{7FA28F}
\definecolor{CalPolyMint}{HTML}{DDF1E7}
\definecolor{CalPolySand}{HTML}{A29F86}
\definecolor{CalPolyBeige}{HTML}{EFEEDE}
\definecolor{DeckInk}{HTML}{1F2933}
\definecolor{DeckMuted}{HTML}{5F6B7A}
\definecolor{DeckSoft}{HTML}{F7F7F5}

\setbeamersize{text margin left=0.55in,text margin right=0.55in}
\setbeamercolor{normal text}{fg=DeckInk,bg=white}
\setbeamercolor{structure}{fg=CalPolyGreen}
\setbeamercolor{title}{fg=white}
\setbeamercolor{frametitle}{fg=CalPolyGreen,bg=white}
\setbeamercolor{block title}{bg=CalPolySage,fg=white}
\setbeamercolor{block body}{bg=CalPolyMint,fg=black}
\setbeamercolor{alerted text}{fg=CalPolyGold}

\setbeamerfont{title}{size=\Huge,series=\bfseries}
\setbeamerfont{frametitle}{size=\Large,series=\bfseries}
\setbeamerfont{block title}{size=\normalsize,series=\bfseries}

\setbeamertemplate{navigation symbols}{}
\setbeamertemplate{itemize item}{\color{CalPolyGold}$\blacktriangleright$}
\setbeamertemplate{itemize subitem}{\color{CalPolyGreen}$\bullet$}
\setbeamertemplate{footline}{
  \hfill{\color{DeckMuted}\scriptsize\insertframenumber/\inserttotalframenumber}
  \hspace{0.18in}\vspace{0.12in}
}
\setbeamertemplate{frametitle}{
  \vspace{0.18in}
  \hspace{0.02in}\insertframetitle\par
  \vspace{0.08in}
  \noindent\makebox[\linewidth][l]{%
    {\color{CalPolyGold}\rule{0.55in}{2pt}}%
    {\color{CalPolyGreen}\rule{\dimexpr\linewidth-0.55in\relax}{2pt}}%
  }
}

\newcommand{\CourseNumber}{Course}
\newcommand{\CourseTitle}{Course Title}
\newcommand{\LectureNumber}{0}
\newcommand{\LectureTitle}{Lecture Title}
\newcommand{\LectureDate}{Date}
\newcommand{\CourseUnit}{Unit}

\newcommand{\coursenumber}[1]{\renewcommand{\CourseNumber}{#1}}
\newcommand{\coursetitle}[1]{\renewcommand{\CourseTitle}{#1}}
\newcommand{\lecturenumber}[1]{\renewcommand{\LectureNumber}{#1}}
\newcommand{\lecturetitle}[1]{\renewcommand{\LectureTitle}{#1}}
\newcommand{\lecturedate}[1]{\renewcommand{\LectureDate}{#1}}
\newcommand{\courseunit}[1]{\renewcommand{\CourseUnit}{#1}}

\author{Kirk Duran}
\institute{California Polytechnic State University, San Luis Obispo}
\date{\LectureDate}

\newcommand{\makedecktitle}{
  \begingroup
  \setbeamercolor{background canvas}{bg=CalPolyGreen}
  \begin{frame}[plain]
    \vfill
    \begin{center}
      {\usebeamerfont{title}\color{white}\LectureTitle\par}
      \vspace{0.18in}
      {\Large\color{white}Lecture \LectureNumber\par}
    \end{center}
    \vfill
    \noindent{\color{white}\rule{\linewidth}{1.2pt}}\par
    \vspace{0.08in}
    \noindent\colorbox{CalPolyGold}{\makebox[0.24\linewidth][c]{\color{white}\Large\CourseNumber}}
    \hspace{0.12in}{\color{white}\Large Kirk Duran}
  \end{frame}
  \endgroup
}

\newenvironment{keyidea}{\begin{block}{Key idea}}{\end{block}}
\newenvironment{activity}{\begin{block}{In class}}{\end{block}}
\newenvironment{cpdefinition}{\begin{block}{Definition}}{\end{block}}
\newenvironment{exampleblockcp}{
  \begingroup
  \setbeamercolor{block title}{bg=CalPolySand,fg=white}
  \setbeamercolor{block body}{bg=CalPolyBeige,fg=black}
  \begin{block}{Example}
}{
  \end{block}
  \endgroup
}

\newcommand{\imageplaceholder}[2][0.3in]{%
  \noindent\makebox[\linewidth][c]{%
    \fcolorbox{CalPolySage}{CalPolyBeige}{%
      \parbox[c][#1][c]{0.86\linewidth}{%
        \centering
        {\color{DeckMuted}\scriptsize Image placeholder: }%
        {\color{DeckInk}\scriptsize\ttfamily\detokenize{#2}}%
      }%
    }%
  }%
}

\newcommand{\sectionframe}[1]{
  \begingroup
  \setbeamercolor{background canvas}{bg=CalPolyGreen}
  \begin{frame}[plain]
    \vfill
    {\color{white}\LARGE\bfseries #1\par}
    \vspace{0.18in}
    {\color{CalPolyGold}\rule{0.22\paperwidth}{3pt}}
    {\color{white}\rule{0.62\paperwidth}{3pt}}
    \vfill
  \end{frame}
  \endgroup
}

\newcommand{\placeholdernote}{
  \begin{block}{Placeholder}
    Replace these bullets with the final lecture material, examples, and in-class activities.
  \end{block}
}


\usepackage{tikz}

\graphicspath{{../assets/lecture-14/}}

\newcommand{\lectureimage}[2][1.75in]{%
  \IfFileExists{../assets/lecture-14/#2}{%
    \includegraphics[width=\linewidth,height=#1,keepaspectratio]{#2}%
  }{}%
}

\newcommand{\lectureplaceholder}[2][2.35in]{%
  \begin{center}
    \fbox{%
      \begin{minipage}[c][#1][c]{0.90\linewidth}
        \centering
        \small\textit{#2}
      \end{minipage}%
    }
  \end{center}%
}

\setbeamersize{text margin left=0.42in,text margin right=0.42in}

\coursenumber{CS 4880}
\coursetitle{Artificial Intelligence}
\lecturenumber{14}
\lecturetitle{Propositional Inference: Logical Equivalence, Truth Tables, and Model Checking}
\lecturedate{September 28, 2026}
\courseunit{Logic and Knowledge Representation}

\begin{document}

% Estimated pacing: 50 minutes
% Entailment review + semantic categories: 7 min
% Truth tables + equivalence: 13 min
% Model checking worked example: 12 min
% Inference rules + soundness/completeness: 10 min
% Unsatisfiability bridge to resolution: 6 min
% Wrap: 2 min

\makedecktitle

\section{From Entailment to Inference}

\begin{frame}[shrink=12]{Last Time: What Does Entailment Mean?}
  \begin{columns}[T,onlytextwidth]
    \begin{column}{0.53\textwidth}
      A knowledge base constrains the set of worlds the agent considers possible.

      \[
        KB \models \alpha
      \]

      means that \textbf{every model satisfying the knowledge base also satisfies $\alpha$}.

      \begin{block}{The missing piece}
        The definition tells us what entailment \emph{means}, but not yet how an agent can decide it computationally.
      \end{block}
    \end{column}
    \begin{column}{0.41\textwidth}
      \lectureimage[2.70in]{slide2.png}
    \end{column}
  \end{columns}
  \begin{keyidea}
    Today we turn the semantic definition of entailment into an inference procedure.
  \end{keyidea}
\end{frame}

\begin{frame}[shrink=12]{The Logical Agent Has a Query}
  \begin{columns}[T,onlytextwidth]
    \begin{column}{0.54\textwidth}
      Suppose a delivery robot knows
      \[
        B \rightarrow C
      \]
      \[
        B
      \]

      where $B$ means ``battery is low'' and $C$ means ``robot should charge.''

      The agent asks:
      \[
        KB \models C\; ?
      \]

      We need a procedure that either confirms the query or finds a counterexample world.
    \end{column}
    \begin{column}{0.40\textwidth}
      \lectureimage[2.70in]{slide3.png}
    \end{column}
  \end{columns}
\end{frame}

\begin{frame}[shrink=12]{Entailment Is a Universal Claim}
  \begin{itemize}
    \item $KB\models\alpha$ says something about \textbf{every} model satisfying $KB$.
    \item To prove entailment semantically, all $KB$-models must make $\alpha$ true.
    \item To disprove entailment, one countermodel is enough.
  \end{itemize}

  \begin{keyidea}
    Entailment is ``for all surviving worlds''; non-entailment needs only one surviving world where the query is false.
  \end{keyidea}
\end{frame}

\begin{frame}[shrink=12]{Countermodels: The Fastest Way to Break a Claim}
  \begin{columns}[T,onlytextwidth]
    \begin{column}{0.54\textwidth}
      Consider the claim
      \[
        P\lor Q \models P.
      \]

      Try the model
      \[
        P=\mathrm{False},\qquad Q=\mathrm{True}.
      \]

      Then $P\lor Q$ is true while $P$ is false.

      \begin{block}{Conclusion}
        That single model disproves the entailment.
      \end{block}
    \end{column}
    \begin{column}{0.40\textwidth}
      \lectureimage[2.45in]{slide5.png}
    \end{column}
  \end{columns}
\end{frame}

\begin{frame}[shrink=12]{Three Useful Semantic Categories}
  \begin{description}
    \item[Tautology:] a sentence true in every model.
    \item[Contradiction:] a sentence false in every model.
    \item[Satisfiable:] a sentence true in at least one model.
  \end{description}

  \begin{block}{Important distinction}
    A sentence can be satisfiable without being a tautology.
  \end{block}

  \begin{keyidea}
    These are properties of a sentence across its possible models.
  \end{keyidea}
\end{frame}

\section{Truth Tables and Logical Equivalence}

\begin{frame}[shrink=12]{Truth Tables Enumerate Possible Worlds}
  \begin{center}
    \renewcommand{\arraystretch}{1.35}
    \begin{tabular}{c|c}
      \textbf{$P$} & \textbf{$Q$} \\\hline
      T & T \\
      T & F \\
      F & T \\
      F & F
    \end{tabular}
  \end{center}

  With two proposition symbols there are
  \[
    2^2=4
  \]
  complete models.
\end{frame}

\begin{frame}[shrink=12]{Truth Tables Give the Meaning of Connectives}
  \begin{center}
    \renewcommand{\arraystretch}{1.25}
    \begin{tabular}{c|c|c|c|c}
      $P$ & $Q$ & $\neg P$ & $P\land Q$ & $P\lor Q$ \\\hline
      T & T & F & T & T \\
      T & F & F & F & T \\
      F & T & T & F & T \\
      F & F & T & F & F
    \end{tabular}
  \end{center}

  \begin{keyidea}
    A truth table is a complete semantic specification of a propositional connective.
  \end{keyidea}
\end{frame}

\begin{frame}[shrink=12]{Implication Revisited}
  \begin{center}
    \renewcommand{\arraystretch}{1.30}
    \begin{tabular}{c|c|c}
      $P$ & $Q$ & $P\rightarrow Q$ \\\hline
      T & T & T \\
      T & F & F \\
      F & T & T \\
      F & F & T
    \end{tabular}
  \end{center}

  $P\rightarrow Q$ is false only when $P$ is true and $Q$ is false.

  \begin{keyidea}
    Implication is a constraint: the forbidden world is $P=T, Q=F$.
  \end{keyidea}
\end{frame}

\begin{frame}[shrink=12]{Biconditional Means ``Same Truth Value''}
  \begin{center}
    \renewcommand{\arraystretch}{1.30}
    \begin{tabular}{c|c|c}
      $P$ & $Q$ & $P\leftrightarrow Q$ \\\hline
      T & T & T \\
      T & F & F \\
      F & T & F \\
      F & F & T
    \end{tabular}
  \end{center}

  $P\leftrightarrow Q$ is true exactly when $P$ and $Q$ agree.
\end{frame}

\begin{frame}[shrink=12]{Logical Equivalence}
  \begin{cpdefinition}
    Two sentences are \textbf{logically equivalent} when they have the same truth value in every model.
  \end{cpdefinition}

  We write
  \[
    \alpha \equiv \beta.
  \]

  Equivalence allows us to replace one expression with another without changing its meaning.

  \begin{keyidea}
    Logical equivalence is semantic sameness, not merely visual similarity.
  \end{keyidea}
\end{frame}

\begin{frame}[shrink=12]{Example: Eliminate an Implication}
  \begin{center}
    \renewcommand{\arraystretch}{1.25}
    \begin{tabular}{c|c|c|c}
      $P$ & $Q$ & $P\rightarrow Q$ & $\neg P\lor Q$ \\\hline
      T & T & T & T \\
      T & F & F & F \\
      F & T & T & T \\
      F & F & T & T
    \end{tabular}
  \end{center}

  The final two columns match in every row, so
  \[
    P\rightarrow Q \equiv \neg P\lor Q.
  \]
\end{frame}

\begin{frame}[shrink=12]{De Morgan's Laws}
  \begin{columns}[T,onlytextwidth]
    \begin{column}{0.54\textwidth}
      \[
        \neg(P\land Q) \equiv \neg P\lor\neg Q
      \]

      \[
        \neg(P\lor Q) \equiv \neg P\land\neg Q
      \]

      Negating a conjunction changes AND to OR; negating a disjunction changes OR to AND.
    \end{column}
    \begin{column}{0.40\textwidth}
      \lectureimage[2.55in]{slide13.png}
    \end{column}
  \end{columns}
\end{frame}

\begin{frame}[shrink=12]{Other Equivalences Worth Knowing}
  \begin{itemize}
    \item Double negation: $\neg\neg P\equiv P$.
    \item Commutativity: $P\land Q\equiv Q\land P$ and $P\lor Q\equiv Q\lor P$.
    \item Associativity lets repeated ANDs or ORs regroup.
    \item Distributivity lets AND and OR move across one another.
    \item Biconditional:
      \[
        P\leftrightarrow Q \equiv (P\rightarrow Q)\land(Q\rightarrow P).
      \]
  \end{itemize}

  \begin{keyidea}
    These identities become tools for automated rewriting next lecture.
  \end{keyidea}
\end{frame}

\begin{frame}[shrink=12]{Why Equivalence Matters to an AI System}
  \begin{columns}[T,onlytextwidth]
    \begin{column}{0.54\textwidth}
      \begin{itemize}
        \item Different-looking sentences can encode exactly the same information.
        \item A reasoning system can rewrite expressions into forms that are easier to process.
        \item A valid rewrite must preserve the set of satisfying models.
        \item Next lecture, we exploit this idea with normal forms.
      \end{itemize}
    \end{column}
    \begin{column}{0.40\textwidth}
      \lectureimage[2.55in]{slide15.png}
    \end{column}
  \end{columns}
\end{frame}

\section{Truth-Table Model Checking}

\begin{frame}[shrink=12]{Model Checking: Use the Definition Directly}
  \begin{columns}[T,onlytextwidth]
    \begin{column}{0.54\textwidth}
      To test $KB\models\alpha$:
      \begin{enumerate}
        \item Enumerate all possible models.
        \item Ignore models where $KB$ is false.
        \item Inspect $\alpha$ in every model where $KB$ is true.
        \item If $\alpha$ is true in all of them, entailment holds.
      \end{enumerate}
    \end{column}
    \begin{column}{0.40\textwidth}
      \lectureimage[2.70in]{slide16.png}
    \end{column}
  \end{columns}
\end{frame}

\begin{frame}[shrink=12]{Worked Example: Does the KB Entail $Q$?}
  \begin{center}
    \Large
    \[
      KB=(P\rightarrow Q)\land P
    \]
    \[
      \text{Query: }Q
    \]
  \end{center}

  We will evaluate the knowledge base and query across all four models for $P$ and $Q$.
\end{frame}

\begin{frame}[shrink=12]{Step 1: Enumerate the Models}
  \begin{center}
    \renewcommand{\arraystretch}{1.25}
    \begin{tabular}{c|c|c|c}
      $P$ & $Q$ & $P\rightarrow Q$ & $KB=(P\rightarrow Q)\land P$ \\\hline
      T & T & T & T \\
      T & F & F & F \\
      F & T & T & F \\
      F & F & T & F
    \end{tabular}
  \end{center}

  Only rows where the $KB$ column is true matter.
\end{frame}

\begin{frame}[shrink=12]{Step 2: Keep Only $KB$-Models}
  \begin{columns}[T,onlytextwidth]
    \begin{column}{0.54\textwidth}
      The knowledge base is true only in the model
      \[
        P=T,\qquad Q=T.
      \]

      Every other model violates at least one sentence in the knowledge base.

      The KB has narrowed four possible worlds down to one.
    \end{column}
    \begin{column}{0.40\textwidth}
      \lectureimage[2.55in]{slide19.png}
    \end{column}
  \end{columns}
\end{frame}

\begin{frame}[shrink=12]{Step 3: Inspect the Query in Every $KB$-Model}
  In the only surviving model, $Q$ is true.

  There is no model where the knowledge base is true and $Q$ is false.

  Therefore
  \[
    KB\models Q.
  \]

  \begin{keyidea}
    The conclusion follows because the knowledge base leaves no possible world in which the query is false.
  \end{keyidea}
\end{frame}

\begin{frame}[shrink=12]{A Non-Entailment Example}
  \begin{columns}[T,onlytextwidth]
    \begin{column}{0.54\textwidth}
      Let
      \[
        KB=P\lor Q
      \]
      with query $P$.

      The model
      \[
        P=F,\qquad Q=T
      \]
      satisfies the knowledge base but falsifies the query.

      Therefore
      \[
        P\lor Q \not\models P.
      \]
    \end{column}
    \begin{column}{0.40\textwidth}
      \lectureimage[2.45in]{slide21.png}
    \end{column}
  \end{columns}
\end{frame}

\begin{frame}[shrink=12]{Truth-Table Entailment as an Algorithm}
  \begin{block}{TT-ENTAILS?$(KB,\alpha)$}
    \ttfamily\small
    for every complete truth assignment $m$:\\[0.04in]
    \quad if $m$ satisfies $KB$:\\
    \qquad if $m$ does not satisfy $\alpha$: return false\\[0.04in]
    return true
  \end{block}

  \begin{keyidea}
    This is exhaustive model checking: search all possible worlds for a countermodel.
  \end{keyidea}
\end{frame}

\begin{frame}[shrink=12]{The Cost: $2^n$ Models}
  \begin{columns}[T,onlytextwidth]
    \begin{column}{0.54\textwidth}
      With $n$ proposition symbols, there are
      \[
        2^n
      \]
      complete truth assignments.

      \begin{itemize}
        \item $10$ symbols $\rightarrow 1{,}024$ models.
        \item $20$ symbols $\rightarrow 1{,}048{,}576$ models.
        \item $30$ symbols $\rightarrow$ more than one billion models.
      \end{itemize}
    \end{column}
    \begin{column}{0.40\textwidth}
      \lectureimage[2.60in]{slide23.png}
    \end{column}
  \end{columns}

  \begin{keyidea}
    The semantics are simple; exhaustive inference scales badly.
  \end{keyidea}
\end{frame}

\begin{frame}[shrink=12]{This Is Another Search Problem}
  \begin{columns}[T,onlytextwidth]
    \begin{column}{0.54\textwidth}
      \begin{itemize}
        \item \textbf{State:} a partial or complete truth assignment.
        \item \textbf{Branching choice:} assign the next proposition true or false.
        \item \textbf{Target counterexample:} a model satisfying $KB$ and falsifying the query.
        \item Truth-table checking explores the complete binary assignment tree.
      \end{itemize}
    \end{column}
    \begin{column}{0.40\textwidth}
      \lectureimage[2.60in]{slide24.png}
    \end{column}
  \end{columns}

  \begin{keyidea}
    A recurring AI theme: reasoning can often be understood as search over a structured space.
  \end{keyidea}
\end{frame}

\section{Inference Procedures}

\begin{frame}[shrink=12]{Entailment vs. Inference}
  \begin{columns}[T,onlytextwidth]
    \begin{column}{0.47\textwidth}
      \begin{block}{Semantic relation}
        \[
          KB\models\alpha
        \]
        $\alpha$ is true in every model of $KB$.
      \end{block}
    \end{column}
    \begin{column}{0.47\textwidth}
      \begin{block}{Syntactic derivation}
        \[
          KB\vdash_i\alpha
        \]
        Inference procedure $i$ derives $\alpha$ from $KB$.
      \end{block}
    \end{column}
  \end{columns}

  \begin{keyidea}
    Do not confuse ``it follows'' with ``this particular procedure found it.''
  \end{keyidea}
\end{frame}

\begin{frame}[shrink=12]{Inference Rules: Reason Without Enumerating Every World}
  An inference rule licenses a new sentence from existing sentences.

  \begin{block}{Example: Modus Ponens}
    \[
      \frac{P\qquad P\rightarrow Q}{Q}
    \]
  \end{block}

  This is a syntactic operation, but we want it to preserve semantic truth.

  \lectureimage[1.35in]{slide26.png}
\end{frame}

\begin{frame}[shrink=12]{Modus Ponens Is Sound}
  \[
    P
  \]
  \[
    P\rightarrow Q
  \]
  \[
    \therefore Q
  \]

  If both premises are true, there is no model in which the conclusion is false.

  \begin{keyidea}
    Modus Ponens never manufactures a false conclusion from true premises.
  \end{keyidea}
\end{frame}

\begin{frame}[shrink=12]{A Tempting but Invalid Rule}
  \begin{columns}[T,onlytextwidth]
    \begin{column}{0.54\textwidth}
      Suppose
      \[
        P\rightarrow Q
      \]
      and
      \[
        Q.
      \]

      Can we conclude $P$? \textbf{No.}

      Wet ground does not imply rain; a sprinkler may also explain it.

      This is the fallacy of \textbf{affirming the consequent}.
    \end{column}
    \begin{column}{0.40\textwidth}
      \lectureimage[2.45in]{slide28.png}
    \end{column}
  \end{columns}
\end{frame}

\begin{frame}[shrink=12]{Soundness}
  \begin{cpdefinition}
    An inference procedure is \textbf{sound} if every sentence it derives is actually entailed.
  \end{cpdefinition}

  \[
    KB\vdash_i\alpha \quad\Rightarrow\quad KB\models\alpha
  \]

  A sound procedure may fail to find some truths, but it will not derive unsupported conclusions.

  \begin{keyidea}
    Soundness protects correctness.
  \end{keyidea}
\end{frame}

\begin{frame}[shrink=12]{Completeness}
  \begin{cpdefinition}
    An inference procedure is \textbf{complete} if it can derive every entailed sentence.
  \end{cpdefinition}

  \[
    KB\models\alpha \quad\Rightarrow\quad KB\vdash_i\alpha
  \]

  A complete procedure cannot permanently miss a conclusion that really follows.

  \begin{keyidea}
    Completeness protects coverage.
  \end{keyidea}
\end{frame}

\begin{frame}[shrink=12]{Truth-Table Model Checking Is Sound and Complete}
  \begin{itemize}
    \item It checks the semantic definition directly.
    \item If it says ``entailed,'' there is no countermodel.
    \item If entailment holds, exhaustive enumeration will confirm it.
    \item The problem is not correctness---it is cost.
  \end{itemize}

  \begin{keyidea}
    We now have a correct baseline algorithm. The next problem is doing less work.
  \end{keyidea}
\end{frame}

\section{Entailment as Unsatisfiability}

\begin{frame}[shrink=12]{Satisfiability Gives Us Another View of Entailment}
  $KB\models\alpha$ exactly when there is no model satisfying $KB$ while $\alpha$ is false.

  That means
  \[
    KB\land\neg\alpha
  \]
  must be unsatisfiable.

  Therefore
  \[
    \boxed{KB\models\alpha \iff KB\land\neg\alpha\text{ is unsatisfiable}.}
  \]

  \begin{keyidea}
    To prove $\alpha$, assume $\neg\alpha$ and show that no consistent world remains.
  \end{keyidea}
\end{frame}

\begin{frame}[shrink=12]{Example: Turn Entailment into a Contradiction Test}
  Want to prove
  \[
    \{P,\;P\rightarrow Q\}\models Q.
  \]

  Add the negation of the query:
  \[
    \neg Q.
  \]

  Now ask whether
  \[
    P\land(P\rightarrow Q)\land\neg Q
  \]
  can be satisfied.

  It cannot: $P$ and $P\rightarrow Q$ force $Q$, while $\neg Q$ forces $Q$ to be false.
\end{frame}

\begin{frame}[shrink=12]{Why This Reformulation Matters}
  \begin{columns}[T,onlytextwidth]
    \begin{column}{0.54\textwidth}
      \begin{itemize}
        \item Truth tables prove entailment by checking all possible worlds.
        \item Unsatisfiability reframes the same problem as deriving an inconsistency.
        \item Perhaps we can manipulate sentences directly instead of explicitly enumerating every model.
      \end{itemize}
    \end{column}
    \begin{column}{0.40\textwidth}
      \lectureimage[2.55in]{slide34.png}
    \end{column}
  \end{columns}
\end{frame}

\begin{frame}[shrink=12]{Next: Automated Logical Inference}
  \begin{columns}[T,onlytextwidth]
    \begin{column}{0.54\textwidth}
      Next lecture we will
      \begin{itemize}
        \item convert sentences into a standardized form;
        \item use resolution to derive contradictions mechanically;
        \item exploit Horn-clause structure;
        \item compare forward and backward chaining.
      \end{itemize}
    \end{column}
    \begin{column}{0.40\textwidth}
      \lectureimage[2.55in]{slide35.png}
    \end{column}
  \end{columns}
\end{frame}

\begin{frame}[shrink=12]{Takeaways}
  \begin{itemize}
    \item Truth tables enumerate models and can decide propositional entailment.
    \item Logical equivalence means identical truth values in every model.
    \item A single countermodel disproves entailment.
    \item Soundness protects correctness; completeness protects coverage.
    \item $KB\models\alpha$ iff $KB\land\neg\alpha$ is unsatisfiable.
    \item Next lecture: reason by manipulating sentences rather than enumerating worlds.
  \end{itemize}

  \begin{keyidea}
    Semantic inference gives us a correct baseline; automated proof procedures aim to keep the guarantees while doing less work.
  \end{keyidea}
\end{frame}

\end{document}
